\section{Which alternating-group subset fields admit the argument?}
\label{sec:subsets}

Let $L/\mathbb Q$ be a Galois extension with group $A_n$, $n\ge4$. Fix
a $k$-element subset of the $n$ letters, and let $K$ be the fixed field
of its setwise stabilizer. Then $[K:\mathbb Q]=\binom nk$. Complementation
allows us to take $1\le k\le n/2$. The following criterion gives an infinite
sequence of even degrees in which the argument of the preceding section works
uniformly in $L$. For a permutation, orbit length here refers to its action
on the set of all $k$-subsets.

\paragraph{Subset criterion.}
Every element of $A_n$ has an odd orbit if and only if either
$\binom nk$ is odd, or
\[
       n\in\{2^a,2^a+1\},\qquad k=2^{a-1},\qquad a\ge2.
\]
In the displayed cases, the maximum over $g\in A_n$ of the length of its
shortest odd orbit is exactly the largest odd integer $m\le n$.
Consequently every rational prime has an odd residue degree at most $m$
somewhere above it in $K$, including primes ramified in $L$.
The bound of the preceding section applies with $m$, independently of
the discriminant of $K$. If the criterion fails, a positive-density set
of unramified primes has exclusively even residue degrees in $K$.

\paragraph{Proof of the criterion.}
A permutation $g$ has an odd orbit on subsets precisely when the Sylow
$2$-subgroup of $\langle g\rangle$ fixes a subset. This follows by taking
the index of its stabilizer in the cyclic group. It suffices, therefore, to
test all even permutations of $2$-power order. Their cycle lengths form
partitions of $n$ into powers of two, with an even number of even parts.
The fixed subset sizes are exactly the subset sums of the parts.

Write $B$ for the powers that occur in the binary expansion of $n$.
Every partition into powers of two refines $B$, as one sees by merging
equal parts until all parts are distinct. Splitting a part leaves every
previously attainable subset size attainable. Therefore every submask of
$n$ (every sum of a subset of $B$) always works. These sizes have
$\binom nk$ odd, by expanding $(1+X)^n=\prod_{2^j\in B}(1+X^{2^j})$
in $\mathbb F_2[X]$.

Suppose $k$ is not a submask. If $B$ itself describes an even permutation,
it excludes $k$. Otherwise, every even partition is a proper refinement,
and replacing any one part $2^j>1$ of $B$ by two halves already gives an
even permutation. Thus $k$ must be a subset sum after \emph{every}
possible single split. This is sufficient too, because a refinement starts
with one of those splits. For a given $j$, a new sum must use exactly one
of the two halves. If $2^{j-1}\in B$ already, the two copies used in such
a sum either leave the binary bit at $j-1$ or carry into the vacated
bit $j$. All sums are then old. If $2^{j-1}\notin B$, a new sum has
exactly one bit outside $B$, at $j-1$. Two different choices of $j$
therefore cannot allow the same new sum. The only remaining $n$ have
a single binary part greater than one, with an optional part $1$.
For $n\ge4$ the single split allows precisely the additional sizes
$2^{a-1}$ and, when $n$ is odd, its complement $2^{a-1}+1$.
This proves the equivalence. For an excluded pair Chebotarev realizes
the conjugacy class of the excluding element with positive density.

\paragraph{The sharp orbit length.}
Put $k=2^{a-1}$. Write each cycle length of $g$ as $2^bu$ with $u$ odd.
Its $2$-part $2^b$ divides $k$ in both listed degrees. Indeed a cycle
violating this would have length $2^a$ and at most one point outside it,
making the permutation odd. Sort the cycles by decreasing $2^b$. Take
whole cycles until the next cycle would reach or exceed $k$ letters.
The remaining number of letters required is a multiple of the $2^b$ of
this cycle. Its $u$th power splits the cycle into $u$ cycles of length
$2^b$. Take the required number of these. The resulting $k$-set is fixed
by $g^u$, so its orbit has odd length dividing $u\le m$.

When $n=2k$ a cycle of length $m=2k-1$ with a fixed letter attains $m$.
A $k$-set selects $k$ or $k-1$ letters in the long cycle, and both
numbers are coprime to $m$. Its stabilizer in this regular cyclic action
has order dividing the number selected. Hence it is trivial. When
$n=2k+1$ the $n$-cycle attains $m=n$ by the same coprimality argument.
Both witness permutations are even. The inertia-orbit argument in the
preceding section gives the assertion about all rational primes.
\hfill$\square$

\begin{center}
\begin{tabular}{rrrrl}
\toprule
$n$ & $k$ & $[K:\mathbb Q]$ & $m$ & Cycle type attaining $m$\\
\midrule
4 & 2 & 6 & 3 & $3+1$\\
5 & 2 & 10 & 5 & $5$\\
8 & 4 & 70 & 7 & $7+1$\\
9 & 4 & 126 & 9 & $9$\\
16 & 8 & 12,870 & 15 & $15+1$\\
17 & 8 & 24,310 & 17 & $17$\\
\bottomrule
\end{tabular}
\end{center}
For example, a degree-$70$ field in the table can use the same universal
odd-residue-degree test as a degree-seven field. The sharpness just proved
concerns orbit lengths. Determining which primes occur as actual isogeny
degrees requires information about elliptic curves. The script
\texttt{checks/alternating.py} enumerates cycle partitions through $n=34$
and independently visits subset orbits of individual permutations through
$n=7$.
